\section{Conclusion}
\label{sec:conclude}

In this paper we provided improved upper and lower bounds on the problem of estimating and sketching effective resistances on expanders. On the algorithmic side we show how sketches tailored to $\ell_1$ when carefully applied to asymmetric formulations of the quadratic form of the Laplacian pseudoinverse gave our results. On the lower bound side, we provided an alternative to the trace estimation approach of \cite{musco2017spectrum} for showing lower bounds and coupled it with techniques of randomly signing edges of the graph to obtain our results. Further, we showed that these techniques had broader implications for addressing algorithmic challenges in numerical linear algebra. 

Beyond the natural open problem of improving both our upper and lower bounds towards bringing them together there are interesting open problems in broadening the applicability of both our upper and lower bounds. For example, obtaining an $\Otilde{m \epsilon^{-1}}$ time algorithm for estimating the effective resistance of all edges in a general (non-expander) graph and extending our $\widetilde{\Omega}(n^2 \epsilon^{-1/2})$ lower bounds to deterministic algorithms remain interesting open problems. We hope that the results of this paper provide useful tools for addressing each. 

